\documentclass[11pt]{article}
\usepackage[T1]{fontenc}
\usepackage{lmodern}
\usepackage[margin=1.1in]{geometry}
\usepackage{amsmath,amssymb,amsthm}

\newtheorem{theorem}{Theorem}[section]
\newtheorem{lemma}[theorem]{Lemma}
\newtheorem{proposition}[theorem]{Proposition}
\newtheorem{corollary}[theorem]{Corollary}
\theoremstyle{definition}
\newtheorem{definition}[theorem]{Definition}
\newtheorem{example}[theorem]{Example}
\theoremstyle{remark}
\newtheorem*{remark}{Remark}

\newcommand{\Z}{\mathbb{Z}}
\DeclareMathOperator{\ord}{ord}

\title{Algebra I --- Lecture 7\\ Subgroups and Cosets}
\author{Lecture notes}
\date{}

\begin{document}
\maketitle

\section{Subgroups}

Throughout, $G$ denotes a group with identity $e$.

\begin{definition}
A nonempty subset $H \subseteq G$ is a \emph{subgroup}, written $H \le G$,
if $ab^{-1} \in H$ for all $a, b \in H$.
\end{definition}

\begin{example}
For every $n \ge 1$ the multiples $n\Z$ form a subgroup of $(\Z,+)$.
\end{example}

\begin{proposition}\label{prop:intersection}
The intersection of any family of subgroups of $G$ is a subgroup of $G$.
\end{proposition}

\begin{proof}
Each subgroup contains $e$, so the intersection is nonempty. If $a$ and $b$
lie in every member of the family, so does $ab^{-1}$.
\end{proof}

\begin{remark}
The union of two subgroups is rarely a subgroup: $2\Z \cup 3\Z$ contains $2$
and $3$ but not $5$.
\end{remark}

\section{Cosets and Lagrange's theorem}

\begin{definition}
For $H \le G$ and $g \in G$, the \emph{left coset} of $H$ containing $g$ is
$gH = \{ gh : h \in H \}$. The number of left cosets is the \emph{index}
$[G : H]$.
\end{definition}

\begin{lemma}\label{lem:partition}
The left cosets of $H$ partition $G$, and any two of them have the same
cardinality.
\end{lemma}

\begin{proof}
The relation $a \sim b \iff a^{-1}b \in H$ is an equivalence relation whose
classes are exactly the left cosets. The map $h \mapsto gh$ is a bijection
from $H$ onto $gH$.
\end{proof}

\begin{theorem}[Lagrange]\label{thm:lagrange}
If $G$ is finite and $H \le G$, then $|G| = [G : H]\,|H|$. In particular
$|H|$ divides $|G|$.
\end{theorem}

\begin{proof}
By Lemma~\ref{lem:partition}, $G$ is the disjoint union of $[G:H]$ cosets,
each with $|H|$ elements.
\end{proof}

\begin{corollary}
For every $g$ in a finite group $G$, $\ord(g)$ divides $|G|$, and
$g^{|G|} = e$.
\end{corollary}

\begin{proof}
Apply Theorem~\ref{thm:lagrange} to the cyclic subgroup
$\langle g \rangle$, whose order is $\ord(g)$.
\end{proof}

\begin{corollary}[Fermat]
If $p$ is prime and $p \nmid a$, then $a^{p-1} \equiv 1 \pmod p$.
\end{corollary}

\begin{example}
The group $S_3$ has order $6$, so its subgroups have order $1$, $2$, $3$ or
$6$. Every one of these orders occurs.
\end{example}

\end{document}
